\begin{helper}
Fix a history cell \(H\), a finite terminal coordinate set \(U\), one blue
support \(B\), and one red label \(J\).  Assume \(0\le p\le1/2\), and assume
the complete terminal product law over \(H\): for every \(A\subseteq U\),
the conditional mass of the complete assignment with present set \(A\) is
\[
p^{|A|}(1-p)^{|U|-|A|}.
\]
Let \(\mathcal T\) be the finite branch-label type and let \(\mathcal P\) be
the finite transcript family.  For a branch label \(\beta\), let
\(\operatorname{blueTrace}(\beta)\subseteq U\) be its complete terminal blue
truth table, and let \(R_{\beta,J}\) be the red support decoded from \(J\)
after the branch is fixed.  Let
\(\mathsf{Realized}(\beta,\tau)\) mean that the pair is a nonempty fixed
\((B,J)\) cell, where \(\tau=(P_\tau,A_\tau,Z_\tau)\).  Let
\(\mathsf{Comp}(A,\beta,\tau)\) be compatibility between a complete terminal
assignment \(A\subseteq U\) and a realized cell.

Assume \(B\) is a blue \(u_B\)-subset of \(U\).  For every realized pair
\((\beta,\tau)\), assume
\[
|R_{\beta,J}|=u_R,\qquad R_{\beta,J}\subseteq U,
\]
that \(R_{\beta,J}\) is red-only, and that
\[
P_\tau,A_\tau\subseteq U,\qquad P_\tau\cap A_\tau=\emptyset,
\]
\[
B\cup R_{\beta,J}\subseteq P_\tau,\qquad Z_\tau\subseteq A_\tau,\qquad
\max(0,a-u_R-u_B)\le |Z_\tau|.
\]
Assume also that, on blue terminal coordinates,
\[
e\in\operatorname{blueTrace}(\beta)\Longleftrightarrow e\in P_\tau,
\qquad
e\notin\operatorname{blueTrace}(\beta)\Longleftrightarrow e\in A_\tau .
\]
Assume branch labels are functional from the complete terminal blue trace:
two labels with the same truth table on all blue coordinates of \(U\) are
equal.  Thus
\(\noderef{FixedSetHistoryCellRedSupportDeterminedByBlueTrace}\) makes the
fixed-\(J\) red support a function of that blue trace, rather than an
independent red-coordinate choice.

Assume the forward compatibility condition: if
\(\mathsf{Comp}(A,\beta,\tau)\), then all coordinates of \(P_\tau\) are in
\(A\), all coordinates of \(A_\tau\) are outside \(A\), and the blue part of
\(A\) is exactly \(\operatorname{blueTrace}(\beta)\).  Assume conversely
that every complete assignment \(A\subseteq U\) satisfying those three
cylinder conditions is compatible with the realized cell.  Assume the
functional scan/disjointness condition: if the same complete assignment is
compatible with two realized pairs, then the two branch labels and the two
transcripts are equal.

Let \(w_{\beta,\tau}\) be cell weights, nonnegative on formal pairs with
\(\tau\in\mathcal P\), zero on non-realized pairs, and satisfying
\[
w_{\beta,\tau}\le p^{|P_\tau|}(1-p)^{|A_\tau|}
\qquad(\tau\in\mathcal P).
\]
In addition, assume the actual residual product-tree normalization for these
fixed \(B,J\) cells.  Namely, there is a finite residual leaf set
\(\Lambda\), a leaf assignment
\[
\iota(\beta,\tau)\in\Lambda,
\]
and a leaf mass \(\nu:\Lambda\to\mathbb R\) such that
\[
0\le\nu(\lambda)\quad(\lambda\in\Lambda),\qquad
\sum_{\lambda\in\Lambda}\nu(\lambda)\le1,
\]
the leaf assignment is injective on realized cells, and for every realized
\((\beta,\tau)\),
\[
p^{|P_\tau|}(1-p)^{|A_\tau|}
\le
p^{u_R}p^{u_B}
(1-p)^{\max(0,a-u_R-u_B)}\,
\nu(\iota(\beta,\tau)).
\]
Here the last displayed inequality is the termwise residual domination: all
mandatory red and blue present factors, and the common selected-absence
factor, have already been exposed before the residual leaf mass is read.

Then there are nonnegative compatible completion block masses
\(\mu_{\beta,\tau}\), indexed by realized \((B,J)\) cells, such that
\[
\sum_{\beta\in\mathcal T}
  \sum_{\tau\in\mathcal P:\mathsf{Realized}(\beta,\tau)}
    \mu_{\beta,\tau}\le1
\]
and, for every realized pair \((\beta,\tau)\),
\[
w_{\beta,\tau}\le
p^{u_R}p^{u_B}
(1-p)^{\max(0,a-u_R-u_B)}\mu_{\beta,\tau}.
\]
\end{helper}
\begin{proof}
Let
\[
L=\max(0,a-u_R-u_B),\qquad
M=p^{u_R}p^{u_B}(1-p)^L .
\]
The complete terminal product law, the compatibility clauses, and the
functional scan condition describe the terminal cylinders represented by the
realized cells.  The additional residual-tree hypotheses state the
normalization that is not a consequence of ordinary disjointness after
dividing by \(M\): realized cells are assigned injectively to residual leaves
whose total mass is at most \(1\), and each full cell cylinder is dominated
by \(M\) times its assigned leaf mass.

Apply \(\noderef{FixedSetHistoryCellRedResidualTreeNormalization}\) to the
fixed \(B,J\) residual tree.  Its hypotheses are exactly the blue-support
data, the realized-cell geometry, the nonnegative normalized leaf masses, and
the injectivity of the leaf assignment on realized cells.  We obtain a
cell-indexed residual mass function \(r_{\beta,\tau}\) satisfying
\[
\sum_{\beta\in\mathcal T}
  \sum_{\tau\in\mathcal P:\mathsf{Realized}(\beta,\tau)}
    r_{\beta,\tau}\le1,
\]
\[
0\le r_{\beta,\tau}
\]
on realized cells, and
\[
r_{\beta,\tau}=\nu(\iota(\beta,\tau))
\]
on realized cells.

Define
\[
\mu_{\beta,\tau}=r_{\beta,\tau}.
\]
The displayed subprobability bound for \(r\) is the required bound for
\(\mu\), and the nonnegativity of \(r\) gives the required nonnegativity of
\(\mu\).

It remains only to check the termwise domination.  For a realized pair
\((\beta,\tau)\), the cylinder estimate for the cell weight gives
\[
w_{\beta,\tau}\le p^{|P_\tau|}(1-p)^{|A_\tau|}.
\]
The residual-tree termwise domination gives
\[
p^{|P_\tau|}(1-p)^{|A_\tau|}
\le
M\,\nu(\iota(\beta,\tau)).
\]
Since \(r_{\beta,\tau}=\nu(\iota(\beta,\tau))\) on realized cells and
\(\mu=r\), this becomes
\[
w_{\beta,\tau}\le M\,\mu_{\beta,\tau}.
\]
Substituting the definition of \(M\) is exactly the displayed bound in the
statement.
\end{proof}
